\section{第6課　吉田さんは　来月　中国へ　行きます}

\subsection{基本课文}

\begin{enumerate}
    \item 吉田さんは　来月　中国\hl{へ}　行きます。
    \item 李さんは　先月　北京\hl{から}　来ました。
    \item 小野さんは　友達\hl{と}　帰りました。
    \item 森さんは　東京\hl{から}　広島\hl{まで}　新幹線\hl{で}　行きます。
\end{enumerate}

甲：いつ　アメリカへ　行きますか。

乙：１０月に　行きます。

甲：駅へ　何で　行きますか。

乙：自転車で　行きます。

甲：だれと　美術館へ　行きますか。

乙：友達と　行きます。

甲：大阪から　上海まで　飛行機で　行きますか。

乙：いいえ、フェリーで　行きます。


\subsection{语法}

\begin{itemize}
    \item 名[场所]へ　動詞：移动方向或目的地
    \item 名1[场所]から　名2[场所]まで　動詞：移动路径或起点终点
    \item 小野さんは友達と帰りました。
    \item 名詞[交通工具]　で　動詞：交通手段
\end{itemize}

\subsection{表达}

\begin{itemize}
    \item たしか
    \item 家、うち
    \item まっすぐ　帰りました
    \item それ　指前文叙述的事
    \item お疲れ様でした
    \item 自転車
\end{itemize}

\subsection{应用课文: 交通機関}

吉田：李さん、昨日は　何時に　アパートへ　帰りましたか。

李：ええと、たしか　十一時半ごろです。

吉田：何で　帰りましたか。

李：電車です。　渋谷まで　電車で　行きました。駅から　アパートまで　歩いて　帰りました。

吉田：小野さんは？

小野：私も　電車です。　駅から　タクシーで　うちへ　帰りました。

李：何時に　帰りましたか。

小野：十二時ごろです。

\vspace{1em}

小野：　森さん、昨夜は　まっすぐ　帰りましたか。

森：いいえ、課長と　いっしょに　銀座へ　行きました。

李：えっ、銀座ですか。

小野：何時に　うちへ　帰りましたか。

森：夜中の二時です。

李：二時ですか。それは　お疲れ様でした。

\subsection{単語}

\begin{vocablist}[2]
    \vocab{先月}
    \vocab{今月}
    \vocab{来月}
    \vocab{夜中}
    \vocab{昨夜}
    \vocab{コンサート}
    \vocab{クリスマス}
    \vocab{誕生日}
    \vocab{こどもの日}
    \vocab{夏休み}
    \vocab{交通機関}
    \vocab{新幹線}
    \vocab{飛行機}
    \vocab{フェリー}
    \vocab{電車}
    \vocab{バス}
    \vocab{タクシー}
    \vocab{美術館}
    \vocab{アパート}
    \vocab{うち}
    \vocab{プール}
    \vocab{友達}
    \vocab{弟}
    \vocab{行く}
    \vocab{帰る}
    \vocab{来る}
    \vocab{たしか}
    \vocab{まっすぐ}
    \vocab{一緒に}
    \vocab{佐藤}
    \vocab{広島}
    \vocab{京都}
    \vocab{北海道}
    \vocab{箱根}
    \vocab{銀座}
    \vocab{渋谷}
    \vocab{新宿}
    \vocab{お疲れ様でした}
    \vocab{お先に失礼します}
    \vocab{歩いて}
    \vocab{大変}
\end{vocablist}
